Mentoring plan

The mentoring of graduate researchers will consist of three components: 
1) Individual mentoring provided by the PIs; 
2) Participation in the mentoring activities at the Department level; 
3) Participation in mentoring activities at the Institution level.

\textbf{Individual mentoring by the PIs.} 

The graduate researchers funded by this project will be advised or co-advised by PIs. 
PIs have extensive experience in graduate mentoring and successfully placing their mentees in high-quality positions in academia and industry. 
The graduate researchers will join the PIs' research labs and collaborate with the PIs on a variety of research projects related to the proposed research. 
More specifically, the PIs plan to provide the following support and guidance to all graduate students involved in the project. 
1) Career Counseling - 
Regular Meetings: Regular one-on-one meetings to discuss career goals, progress, and strategies.
Networking Opportunities: Facilitation of networking sessions with researchers, industry professionals, and alumni. 
2) Training in Publication and Presentation Preparation - 
Publication Guidance: Mentorship in writing and submitting research papers, as well as writing rebuttals.
Presentation Skills: Training on effective presentation skills, including crafting presenting slides and efficiently conveying information to diverse audiences. 
3) Teaching and Mentoring Skills -
Providing graduate students with opportunities to teach classes and mentor junior researchers (i.e. undergraduate students). 
4) Collaboration Skills - 
Interdisciplinary Collaboration: Encouragement and facilitation of collaborative projects with researchers from various disciplines and institutions. 

Finally, PIs will personally advise graduate researchers on how to search for academic jobs, what to expect at an interview, what to expect in the first few years of the academic career, and other issues related to finding an academic job.

\textbf{Participation in departmental mentoring activities.} 
The graduate researchers will be encouraged to attend the departmental mentoring events, including seminars on effective presentations and writing skills, effective teaching techniques, and career development seminars.

\textbf{Participation in mentoring activities at the institutional level.} 
The PIs will encourage graduate researchers to attend mentoring activities at the institutional level. 
For example, the Center for Science and Engineering Partnerships (CSEP) Professional Development Series (PDS) Program at
UCSB is an intensive Professional Development program offered at UCSB by CSEP. 
The mission of CSEP is to strengthen UCSB’s capacity to play a leading role in the education and professional development of current and future scientists and engineers. 
Activities include efforts to increase scientific literacy in Science, Technology, Engineering, and Mathematics (STEM) fields to meet national workforce needs. 
CSEP was created to provide a central hub on the UCSB campus to develop and administer educational programs in science and engineering and serve as a campus-wide nexus for innovative projects that address state and national needs of broadening participation in STEM, for K-12 through graduate student and postdoctoral scholar levels. 
CSEP, together with multiple groups, departments, faculty, and staff at UCSB, offers an intensive Professional Development Series that focuses on the knowledge and practical experience essential for the success of postdoctoral
scholars and their future varied careers in science, engineering, and mathematics. 
The program covers topics such as Basic skills, Communication skills, and Career skills. 
Workshops and seminars are usually offered bi-weekly, are approximately 1-2 hours in duration, and are led by professionals in the field or faculty from various STEM departments.